\newpage
\chapter{Contact Formulations}
\label{chap:contact-formulations}

Contact describes the interaction of two bodies whose surfaces may touch,
transmit forces and subsequently separate or slide. At first sight this may
look like another application of the \textbf{multifreedom constraints} treated
in the previous chapter. There is, however, one fundamental difference.

An ordinary MFC is an \textit{equality constraint}. Once imposed, it remains
active throughout the analysis. Contact in the normal direction is generally a
\textit{unilateral constraint}: the surfaces must not penetrate each other, but
they are allowed to separate. Consequently, the constraint is active only while
the surfaces are closed.

The simplest scalar distinction is

\begin{equation}
    \underbrace{a(\m{u}) = 0}_{\text{equality constraint}}
    \qquad \text{versus} \qquad
    \underbrace{g_n(\m{u}) \geq 0}_{\text{unilateral contact constraint}}.
\end{equation}

Here $ g_n $ is the \textit{normal gap}. In this chapter the sign convention is:

\begin{equation}
    g_n > 0 \quad \text{open gap}, \qquad
    g_n = 0 \quad \text{closed contact}, \qquad
    g_n < 0 \quad \text{penetration}.
\end{equation}

A contact definition is therefore not merely a special stiffness placed between
two nodes. A complete definition has to specify at least four things:

\begin{enumerate}
    \item how the geometrical gap and relative sliding are measured,
    \item which normal and tangential motions are admissible,
    \item how the corresponding constraints are imposed numerically,
    \item how points on one discretized surface interact with the other surface.
\end{enumerate}

Terms such as \textit{frictionless}, \textit{rough}, \textit{frictional},
\textit{no separation} and \textit{bonded} primarily describe the admissible
relative motion and the tractions that may be transmitted. They should not be
confused with \textit{penalty}, \textit{Lagrange multiplier} and
\textit{augmented Lagrangian}, which describe numerical enforcement methods.
Nor should they be confused with \textit{node-to-surface} or
\textit{surface-to-surface}, which describe contact discretizations.

\section{A One-Dimensional Contact Example}

The basic behaviour of normal contact can be understood without surface
geometry. Consider a node attached to a spring of stiffness $ k $. A positive
external force $ f $ moves the node to the right. A rigid wall is located at
$ u = g_0 $, as sketched conceptually by

\begin{equation}
    \text{spring} \; k \quad \longrightarrow \quad u
    \qquad \vert \qquad \text{rigid wall at } g_0.
\end{equation}

The distance from the node to the wall is

\begin{equation}\label{contact-1d-gap}
    g_n = g_0 - u.
\end{equation}

Thus $ u < g_0 $ means that the node and wall are separated, $ u = g_0 $
means that they touch, and $ u > g_0 $ would represent an inadmissible
penetration.

If the wall is absent, equilibrium gives

\begin{equation}
    k u = f,
\end{equation}

and therefore $ u = f/k $. Two cases are possible.

\begin{enumerate}
    \item If $ f/k < g_0 $, the unconstrained displacement does not reach the
        wall. The contact remains open and the wall exerts no force.

    \item If $ f/k \geq g_0 $, the unconstrained displacement would cross the
        wall. Contact must become active, the displacement is limited to
        $ u = g_0 $, and the wall supplies a reaction force.
\end{enumerate}

Let the compressive contact reaction be denoted by $ p_n \geq 0 $. It acts in
the direction opposite to increasing $ u $. Equilibrium of the node is

\begin{equation}\label{contact-1d-equilibrium}
    k u + p_n = f.
\end{equation}

The exact solution is therefore

\begin{equation}\label{contact-1d-exact-solution}
    \begin{aligned}
        u &= \frac{f}{k},       & p_n &= 0,
        && \text{if } f < k g_0, \\
        u &= g_0,               & p_n &= f-k g_0,
        && \text{if } f \geq k g_0.
    \end{aligned}
\end{equation}

This small example already contains the central difficulty of contact: before
solving the equilibrium equations, it is generally not known whether the
contact reaction is zero or whether the gap is zero. The set of active contact
constraints is part of the solution.

\section{The Unilateral Normal Constraint}

For a frictionless normal contact pair the admissible states are described by
three conditions:

\begin{bbox}
    \begin{equation}\label{contact-kkt-normal}
        g_n \geq 0, \qquad p_n \geq 0, \qquad p_n g_n = 0.
    \end{equation}
\end{bbox}

These are the \textit{Karush--Kuhn--Tucker} conditions for unilateral contact.
Their meaning is more important than their name.

\begin{enumerate}
    \item $ g_n \geq 0 $ states that penetration is not admissible.

    \item $ p_n \geq 0 $ states that ordinary contact can push but cannot pull.
        A tensile normal traction would require an adhesive or bonded interface.

    \item $ p_n g_n = 0 $ is the \textit{complementarity condition}. Either
        the gap is open and the pressure is zero, or the pressure is positive
        and the gap is closed. They cannot both be positive.
\end{enumerate}

The two admissible states may be written explicitly as

\begin{equation}\label{contact-open-closed-states}
    \begin{aligned}
        \text{open:}  && g_n &> 0, & p_n &= 0, \\
        \text{closed:}&& g_n &= 0, & p_n &\geq 0.
    \end{aligned}
\end{equation}

At the instant of first touch both $ g_n = 0 $ and $ p_n = 0 $ are possible.
This transition point separates the open and closed branches.

Equation \eqref{contact-kkt-normal} also explains why replacing contact by a
permanent equality $ g_n = 0 $ is generally incorrect. Such an equality would
prevent separation and could transmit tension. It would describe a
\textit{no-separation} or tied normal interface, not ordinary contact.

\subsection{Contact Status as an Active Set}

Suppose that a finite element model contains many possible contact locations.
At a given nonlinear iteration the solver classifies each location as active or
inactive:

\begin{equation}
    \begin{aligned}
        \mathcal{A} &= \{ i \; | \; \text{contact constraint } i
        \text{ is active} \}, \\
        \mathcal{I} &= \{ i \; | \; \text{contact constraint } i
        \text{ is inactive} \}.
    \end{aligned}
\end{equation}

For an inactive constraint, $ p_{n,i} = 0 $. For an active constraint, the
solver attempts to enforce $ g_{n,i} = 0 $, either exactly or approximately.
After the displacement update, the status is checked again. A contact may close,
open, stick or slip, so the active set changes during the nonlinear solution.

This status update is one reason why a structurally linear model may still
require a nonlinear contact analysis. Even if the material law and the
small-displacement stiffness matrix are linear, the switch between
$ p_n = 0 $ and $ g_n = 0 $ is not.

\section{Contact Kinematics}

The one-dimensional gap \eqref{contact-1d-gap} is replaced in a continuum model
by the distance between two candidate surfaces. It is common to label one side
\textit{slave} and the other side \textit{master}. These labels describe the
search and projection procedure; they do not imply that one physical body is
more important than the other.

Let a point on the slave surface have reference position $ \m{X}_s $ and current
position

\begin{equation}
    \m{x}_s = \m{X}_s + \m{u}_s.
\end{equation}

Let the master surface be parametrized by surface coordinates $ \M{\xi} $:

\begin{equation}
    \m{x}_m(\M{\xi}) = \m{X}_m(\M{\xi}) + \m{u}_m(\M{\xi}).
\end{equation}

A corresponding master point $ \m{x}_m(\closure{\M{\xi}}) $ is found, most often
by a closest-point projection. If $ \m{n}_m $ is the outward unit normal of the
master surface at that point, the normal gap is

\begin{bbox}
    \begin{equation}\label{contact-normal-gap-general}
        g_n = \left( \m{x}_s -
        \m{x}_m(\closure{\M{\xi}}) \right)^T \m{n}_m.
    \end{equation}
\end{bbox}

The orientation of $ \m{n}_m $ must be chosen consistently with the adopted sign
convention. Reversing the normal reverses the sign of $ g_n $ and of the normal
traction definition.

\subsection{Closest-Point Projection}

For a smooth master surface, the closest point satisfies orthogonality of the
slave-to-master vector to both surface tangent directions:

\begin{equation}\label{contact-closest-point-condition}
    \left( \m{x}_s - \m{x}_m(\closure{\M{\xi}}) \right)^T
    \frac{\partial \m{x}_m}{\partial \xi_\alpha} = 0,
    \qquad \alpha = 1,2.
\end{equation}

Equation \eqref{contact-closest-point-condition} is itself nonlinear when the
surface is curved or undergoing finite motion. In practice the projection has
to determine both the contacted element or segment and the local coordinates
inside that element.

If the projection falls outside the candidate face, the solver may test an
adjacent face, edge or vertex. Consequently, contact search is a geometrical
problem in addition to the mechanical equilibrium problem.

\subsection{Normal and Tangential Decomposition}

The normal projector and the tangential projector are

\begin{equation}
    \m{P}_n = \m{n} \m{n}^T,
    \qquad
    \m{P}_t = \m{I} - \m{n} \m{n}^T.
\end{equation}

Any relative displacement increment $ \Delta \m{u}_{rel} $ may be decomposed as

\begin{equation}
    \Delta \m{u}_{rel}
    = \underbrace{\m{P}_n \Delta \m{u}_{rel}}_{\text{normal part}}
    + \underbrace{\m{P}_t \Delta \m{u}_{rel}}_{\text{tangential part}}.
\end{equation}

The incremental tangential slip is commonly written

\begin{equation}\label{contact-tangential-slip-increment}
    \Delta \m{g}_t =
    \m{P}_t \left( \Delta \m{u}_s - \Delta \m{u}_m \right).
\end{equation}

For finite sliding, the tangential basis rotates and the projected master point
moves across element boundaries. Therefore $ \m{g}_t $ is normally updated
incrementally rather than treated as a single difference between initial and
current Cartesian coordinates.

\subsection{Small and Finite Sliding}

A \textit{small-sliding} formulation assumes that the contacted point remains
near its initial projection. The contact normal and the interpolation of the
opposite surface can then be updated only mildly or even kept fixed within an
increment. This is efficient when relative tangential motion is small compared
with the local element dimensions.

A \textit{finite-sliding} formulation allows the slave point to travel over
many master faces. The closest point, surface normal, tangent basis and contact
interpolation must be updated as the bodies move. This is more general but also
more expensive and more sensitive to surface discretization.

Small versus finite sliding is a kinematic choice. It is independent of whether
the interface is frictionless or frictional and independent of whether a
penalty or multiplier method is used.

\section{Normal Constraint Enforcement}

The conditions \eqref{contact-kkt-normal} state the exact mechanical problem,
but they do not by themselves prescribe how the algebraic equations should be
formed. The same three families introduced for MFCs appear again:

\begin{enumerate}
    \item penalty enforcement,
    \item Lagrange multiplier enforcement,
    \item augmented Lagrangian enforcement.
\end{enumerate}

The important difference from ordinary MFCs is that these methods are applied
only to the currently active contact constraints.

\subsection{The Penalty Method}

Define the positive-part or Macaulay bracket by

\begin{equation}
    \langle x \rangle = \max(x,0).
\end{equation}

With the present gap convention, a simple normal penalty law is

\begin{bbox}
    \begin{equation}\label{contact-normal-penalty-law}
        p_n = \epsilon_n \langle -g_n \rangle,
    \end{equation}
\end{bbox}

where $ \epsilon_n > 0 $ is the normal penalty stiffness. If the gap is open,
$ g_n > 0 $, then $ p_n = 0 $. If penetration occurs, $ g_n < 0 $, then the
contact pressure grows in proportion to the penetration $ -g_n $.

The corresponding penalty potential per contact point is

\begin{equation}\label{contact-normal-penalty-potential}
    \Pi_c = \frac{1}{2} \epsilon_n
    \langle -g_n \rangle^2.
\end{equation}

The contact force follows from variation of this potential with respect to the
nodal displacements. If $ \m{B}_n = \partial g_n / \partial \m{u} $ denotes the
gap-displacement matrix, the contact contribution to the residual has the form

\begin{equation}
    \m{r}_c = -\m{B}_n^T p_n,
\end{equation}

up to the sign convention used for the global residual.

\subsubsection{Penalty Solution of the One-Dimensional Example}

For the node-wall example, penetration is $ u-g_0 $. Once contact is active,
choose

\begin{equation}
    p_n = \epsilon_n (u-g_0).
\end{equation}

Substitution into equilibrium \eqref{contact-1d-equilibrium} gives

\begin{equation}
    k u + \epsilon_n (u-g_0) = f,
\end{equation}

and therefore

\begin{equation}\label{contact-1d-penalty-displacement}
    u = \frac{f + \epsilon_n g_0}{k+\epsilon_n}.
\end{equation}

The penetration error is

\begin{equation}\label{contact-1d-penalty-penetration}
    u-g_0 = \frac{f-k g_0}{k+\epsilon_n}.
\end{equation}

As $ \epsilon_n \rightarrow \infty $, the penetration tends to zero and
$ u \rightarrow g_0 $. For any finite $ \epsilon_n $, however, some penetration
is required to generate the contact reaction. This is not a geometrical prediction
of the physical model. It is the constraint violation introduced by penalty
enforcement.

The tradeoff is the same as for penalty MFCs. A small $ \epsilon_n $ gives a
soft interface and excessive penetration. A very large $ \epsilon_n $ improves
constraint accuracy but increases the condition number of the tangent stiffness
matrix and can impair Newton convergence.

\subsection{Choosing the Normal Penalty Stiffness}

There is no universal penalty stiffness that is suitable for every contact pair.
The value has to be large relative to the normal compliance of the adjacent
finite elements, but not so large that the contact terms dominate the structural
stiffness by many unnecessary orders of magnitude.

A useful dimensional estimate is

\begin{equation}\label{contact-penalty-scaling}
    \epsilon_n \sim \alpha \frac{E_{eff} A_c}{h_c},
\end{equation}

for a discrete contact point carrying representative area $ A_c $, or

\begin{equation}
    \bar{\epsilon}_n \sim \alpha \frac{E_{eff}}{h_c}
\end{equation}

for a traction-based surface formulation. Here $ E_{eff} $ is an effective
normal modulus of the contacting materials, $ h_c $ is a representative element
length in the contact-normal direction, and $ \alpha $ is a dimensionless
scaling factor.

These expressions are scaling arguments, not exact laws. Element type,
integration scheme, shell thickness, material incompressibility and the chosen
surface discretization all affect the useful range. The computed penetration
and nonlinear convergence should therefore be checked rather than relying on a
large default number alone.

\subsection{The Lagrange Multiplier Method}

A Lagrange multiplier can enforce the active normal constraint exactly. For a
known active set, collect the active gaps in vector $ \m{g}_n $ and linearize
them as

\begin{equation}
    \Delta \m{g}_n = \m{A}_n \Delta \m{u},
\end{equation}

where $ \m{A}_n $ is the active normal constraint matrix. The Newton correction
has the bordered form

\begin{equation}\label{contact-lagrange-linear-system}
    \begin{bmatrix}
        \m{K}_T & \m{A}_n^T \\
        \m{A}_n & \m{0}
    \end{bmatrix}
    \begin{bmatrix}
        \Delta \m{u} \\
        \Delta \M{\lambda}_n
    \end{bmatrix}
    = -
    \begin{bmatrix}
        \m{r} \\
        \m{g}_n
    \end{bmatrix}.
\end{equation}

The multiplier $ \M{\lambda}_n $ represents the contact pressure or an
equivalent nodal contact force, depending on the discretization. The gap of an
active constraint can be driven to zero without selecting a penalty stiffness.

This exactness has a price. Additional unknowns are introduced, the coefficient
matrix becomes indefinite, and the active set must still be determined. A
multiplier that becomes tensile is inadmissible for ordinary contact; the
corresponding constraint must be released.

For the one-dimensional example, active contact gives the exact equations

\begin{equation}
    \begin{bmatrix}
        k & 1 \\
        1 & 0
    \end{bmatrix}
    \begin{bmatrix}
        u \\
        p_n
    \end{bmatrix}
    =
    \begin{bmatrix}
        f \\
        g_0
    \end{bmatrix}.
\end{equation}

The solution is $ u=g_0 $ and $ p_n=f-k g_0 $. If this computed $ p_n $ is
negative, the assumed active status was wrong and the contact must be opened.

\subsection{The Augmented Lagrangian Method}

The augmented Lagrangian method combines a moderate penalty stiffness with an
iteratively corrected pressure estimate. With iteration index $ j $, one common
normal update is

\begin{bbox}
    \begin{equation}\label{contact-augmented-lagrange-update}
        p_n^{j+1} =
        \left\langle p_n^j - \epsilon_n g_n^{j+1} \right\rangle.
    \end{equation}
\end{bbox}

If penetration remains, $ g_n<0 $, the pressure estimate increases. If the gap
opens, the positive-part projection prevents a tensile contact pressure and
ultimately returns the pressure to zero.

The displacement problem inside an augmentation step resembles a penalty
problem, but the multiplier estimate carries information from one augmentation
to the next. Hence the gap can be reduced without increasing $ \epsilon_n $ to
an excessively large value.

The augmented Lagrangian method is therefore often a useful compromise:

\begin{enumerate}
    \item it is more accurate than a penalty method using the same stiffness,
    \item it avoids the full multiplier system in some implementations,
    \item it requires additional status and augmentation iterations.
\end{enumerate}

The exact implementation varies. Some formulations keep explicit multiplier
unknowns, while others update pressure-like history variables outside the main
Newton solve. The mechanical meaning remains the same: a penalty regularization
is combined with iterative multiplier correction.

\subsection{Assessment of Normal Enforcement Methods}

The three methods should be compared by their numerical consequences rather
than by the apparent elegance of their equations.

\begin{enumerate}
    \item \textit{Penalty enforcement} retains the original displacement unknowns
        and is relatively easy to assemble. The contact constraint is only
        approximate and its accuracy depends on the penalty stiffness.

    \item \textit{Lagrange multiplier enforcement} can satisfy the active gap
        exactly and gives the contact reactions directly. It adds unknowns,
        produces an indefinite system and requires careful treatment of dependent
        constraints.

    \item \textit{Augmented Lagrangian enforcement} reduces the dependence on a
        very large penalty stiffness, but adds an outer correction process and
        still requires contact-status decisions.
\end{enumerate}

No method removes the underlying nonsmoothness caused by opening and closing.
A sophisticated enforcement method cannot compensate for an incorrect contact
normal, poor surface discretization or a failed search projection.

\section{Tangential Contact and Friction}

Normal contact decides whether the surfaces may approach or separate.
Tangential contact decides whether they may slide and which shear traction is
transmitted while they are closed.

Let $ \m{t}_t $ denote the tangential traction vector and $ p_n $ the compressive
normal pressure. A friction law acts only while contact is closed. If the
surfaces are open, both normal and tangential contact tractions vanish:

\begin{equation}
    g_n>0 \quad \Longrightarrow \quad p_n=0, \qquad \m{t}_t=\m{0}.
\end{equation}

\subsection{Frictionless Contact}

Frictionless contact imposes the unilateral normal condition but no tangential
constraint:

\begin{equation}\label{contact-frictionless-definition}
    g_n \geq 0, \qquad p_n \geq 0, \qquad p_n g_n=0,
    \qquad \m{t}_t=\m{0}.
\end{equation}

The surfaces may separate and may slide freely. Compressive normal traction is
transmitted, but no shear traction and no tensile normal traction are carried.

Frictionless does not mean that the bodies are disconnected. While closed, the
normal relative motion is constrained and normal forces can be large. It only
means that the tangential relative motion does not generate contact shear.

\subsection{Coulomb Friction}

The classical Coulomb law limits the tangential traction by the current normal
pressure:

\begin{bbox}
    \begin{equation}\label{contact-coulomb-yield-function}
        \phi = \norm{\m{t}_t} - \mu p_n \leq 0,
    \end{equation}
\end{bbox}

where $ \mu \geq 0 $ is the coefficient of friction. Two states are possible.

\begin{enumerate}
    \item \textit{Stick}: $ \phi < 0 $. The tangential traction lies inside the
        friction bound and the incremental tangential slip is constrained to
        zero, apart from any tangential penalty compliance.

    \item \textit{Slip}: $ \phi = 0 $. The traction magnitude reaches
        $ \mu p_n $ and its direction opposes the sliding direction.
\end{enumerate}

For nonzero sliding velocity $ \m{v}_t $, the ideal slip law is

\begin{equation}\label{contact-coulomb-slip-law}
    \m{t}_t = -\mu p_n
    \frac{\m{v}_t}{\norm{\m{v}_t}}.
\end{equation}

The minus sign states that friction resists relative sliding. At zero velocity
the direction is undefined, which is why the stick state has to be treated as
an inequality rather than by direct substitution into
\eqref{contact-coulomb-slip-law}.

\subsection{A Tangential Trial-State Algorithm}

The stick-slip update resembles an elastic predictor and plastic corrector.
Let $ \epsilon_t $ be a tangential penalty stiffness and let
$ \Delta \m{g}_t $ be the tangential relative-displacement increment. A trial
traction is

\begin{equation}\label{contact-friction-trial-traction}
    \m{t}_t^{trial} = \m{t}_t^n
    - \epsilon_t \Delta \m{g}_t.
\end{equation}

The trial friction function is

\begin{equation}
    \phi^{trial} = \norm{\m{t}_t^{trial}} - \mu p_n.
\end{equation}

If $ \phi^{trial} \leq 0 $, the contact sticks and

\begin{equation}
    \m{t}_t^{n+1} = \m{t}_t^{trial}.
\end{equation}

If $ \phi^{trial} > 0 $, the contact slips and the traction is returned to the
friction limit:

\begin{equation}\label{contact-friction-return}
    \m{t}_t^{n+1} =
    \mu p_n
    \frac{\m{t}_t^{trial}}{\norm{\m{t}_t^{trial}}}.
\end{equation}

The signs in \eqref{contact-friction-trial-traction} depend on the adopted
traction and slip conventions. The important construction is invariant: first
predict the traction assuming stick, then check the Coulomb bound, and if the
bound is exceeded project the traction back to the admissible friction surface.

Because the limit $ \mu p_n $ changes with normal pressure, the normal and
tangential problems are coupled. Opening the contact immediately removes the
friction traction. A change in normal pressure changes the maximum transferable
shear.

\subsection{Regularized Friction}

Ideal Coulomb friction is nonsmooth at the transition from stick to slip and at
zero sliding velocity. Some formulations replace the discontinuous law by a
smooth approximation, for example

\begin{equation}
    \m{t}_t = -\mu p_n
    \tanh\left(\frac{\norm{\m{v}_t}}{v_0}\right)
    \frac{\m{v}_t}{\norm{\m{v}_t}},
\end{equation}

with a small regularization velocity $ v_0 $. Such a law can improve nonlinear
convergence, but it introduces artificial microslip and rate dependence. The
regularization parameter is therefore part of the numerical model and should be
reported when it materially affects the result.

\section{Common Contact Definitions}

The common names used in finite element programs can now be stated in terms of
normal and tangential admissibility. The names are convenient, but their exact
implementation may differ between programs. The solver-independent mechanical
meaning is described below.

\subsection{Frictionless}

A frictionless interface has unilateral normal contact and zero tangential
traction:

\begin{equation}
    g_n \geq 0, \qquad p_n \geq 0, \qquad p_n g_n=0,
    \qquad \m{t}_t=\m{0}.
\end{equation}

It can transmit compression. It cannot transmit tension or shear. It allows
both separation and unrestricted sliding.

\subsection{Frictional}

A frictional interface has unilateral normal contact together with the Coulomb
bound

\begin{equation}
    \norm{\m{t}_t} \leq \mu p_n.
\end{equation}

It can transmit compression and a pressure-dependent amount of shear. It allows
separation. While closed, it may stick or slide depending on whether the required
shear traction is below or at the friction limit.

The frictionless definition is recovered as the special case $ \mu=0 $.

\subsection{Rough}

A rough interface prevents tangential sliding while the surfaces are in contact:

\begin{equation}
    g_n=0 \quad \Longrightarrow \quad \Delta \m{g}_t=\m{0}.
\end{equation}

It may be viewed as the limiting case of friction with an unbounded friction
capacity, although implementations often impose the tangential constraint
directly rather than using a very large value of $ \mu $.

Rough contact may transmit whatever tangential traction is required to prevent
slip while the normal contact remains active. It still permits separation and
normally cannot transmit tensile normal traction. This distinguishes rough
contact from bonded contact.

\subsection{No Separation}

A no-separation interface prevents opening in the normal direction once the
interface is closed, but allows tangential sliding:

\begin{equation}\label{contact-no-separation-definition}
    g_n = 0, \qquad \m{t}_t=\m{0}
\end{equation}

for the ideal frictionless no-separation case. Because the normal relation is
bilateral, the interface can develop both compressive and tensile normal
reactions. It is therefore not ordinary unilateral contact.

The phrase ``once closed'' is important. Some implementations tie surfaces that
are initially within a specified tolerance; others first allow them to close
and then prevent reopening. Initial status and program convention must therefore
be checked.

No separation is useful when two surfaces may slide but are physically retained
against opening, for example by an idealized enclosure. It is not appropriate
when real lift-off is expected.

\subsection{Bonded}

An ideal bonded interface enforces displacement continuity in both normal and
tangential directions:

\begin{bbox}
    \begin{equation}\label{contact-bonded-definition}
        \m{u}_s - \m{u}_m = \m{0}.
    \end{equation}
\end{bbox}

Equivalently,

\begin{equation}
    g_n=0, \qquad \m{g}_t=\m{0}.
\end{equation}

A bonded interface can transmit compression, tension and shear. It allows
neither separation nor sliding. Mechanically it is a tied interface or a set of
MFCs, even if it is created through contact-surface machinery.

Bonded contact should not be interpreted as a universal model of adhesive
failure. The ideal constraint has no strength limit and no fracture energy.
Debonding requires an additional interface law, such as a cohesive-zone model,
a damage law or an explicit release criterion.

\subsection{Comparison of the Definitions}

The principal differences are summarized below. ``Yes'' means that the idealized
definition admits the motion or force without adding another constitutive law.

\begin{center}
\resizebox{\textwidth}{!}{%
\begin{tabular}{|l|c|c|c|c|}
    \hline
    Definition & Separation & Sliding & Tensile normal force & Shear transfer \\
    \hline
    Frictionless & Yes & Yes & No & No \\
    Frictional & Yes & Conditional & No & Limited by $\mu p_n$ \\
    Rough & Yes & No while closed & No & Unlimited while closed \\
    No separation & No & Yes & Yes & No, if frictionless \\
    Bonded & No & No & Yes & Yes \\
    \hline
\end{tabular}%
}
\end{center}

This table describes the ideal mechanics, not the numerical enforcement. For
example, bonded contact may be implemented by direct elimination, MFCs, penalty
terms, multipliers or mortar constraints. Two models with the same word
``bonded'' may therefore have different constraint accuracy and conditioning.

\subsection{Cohesive and Adhesive Interfaces}

A cohesive interface is different from an ideal bond because the transmitted
traction is related to an interface separation. In a simple uncoupled form,

\begin{equation}
    t_n = t_n(\delta_n), \qquad
    \m{t}_t = \m{t}_t(\M{\delta}_t),
\end{equation}

where $ \delta_n $ and $ \M{\delta}_t $ are normal and tangential separations.
The law may initially be stiff, reach a strength limit and then soften as damage
grows. The area under a traction-separation curve represents fracture energy.

Thus a cohesive interface is a constitutive model with finite strength and
possible failure. Bonded contact is an ideal kinematic constraint with unlimited
strength unless an additional failure rule is supplied.

\section{Contact Discretizations}

The continuous contact conditions apply over surfaces, but a finite element
model contains nodes, edges, faces and interpolation functions. A contact
discretization determines where the gap is evaluated and how the resulting
tractions are distributed to the nodes.

\subsection{Node-to-Node Contact}

Node-to-node contact connects predefined pairs of nodes. For a pair with normal
$ \m{n} $, the gap may be written

\begin{equation}
    g_n = (\m{x}_s-\m{x}_m)^T \m{n}.
\end{equation}

The method is simple and inexpensive, but it requires compatible node locations
and a known pairing. It is poorly suited to general sliding because a node that
moves tangentially no longer remains opposite its original partner.

Node-to-node contact is best understood as a collection of discrete contact
springs or constraints, not as an accurate representation of arbitrary smooth
surface interaction.

\subsection{Node-to-Surface Contact}

In node-to-surface contact, each slave node is projected onto a master face.
The gap is evaluated between the slave node and the interpolated master point.
The contact force acting at the projection is distributed to the master nodes
through the face shape functions.

This permits a slave node to move over the master mesh and does not require
matching nodes. However, the result is biased: penetration is prevented at the
slave nodes, but not necessarily at points between them. A coarse slave surface
may therefore pass locally through the master surface between slave nodes.

As a general practical rule, the more finely discretized or more deformable
surface is often selected as slave, while the smoother, coarser or stiffer
surface is selected as master. This is a useful starting point, not a universal
law. Curvature, shell normals, expected sliding direction and element order also
matter.

\subsection{Surface-to-Surface Contact}

Surface-to-surface formulations evaluate contact over slave segments or
integration points rather than only at slave nodes. The contact virtual work may
be written schematically as

\begin{equation}\label{contact-virtual-work}
    \delta W_c =
    \int_{\Gamma_c}
    \left( p_n \delta g_n
    + \m{t}_t^T \delta \m{g}_t \right) d\Gamma.
\end{equation}

Numerical quadrature converts this surface integral into contact contributions
at integration points. Surface-to-surface contact generally represents pressure
more smoothly and reduces nodal bias, especially for higher-order elements and
curved surfaces.

The term does not uniquely identify one algorithm. Different programs may use
one-pass, two-pass, segment-to-segment or mortar-like formulations under the
broad surface-to-surface name.

\subsection{Mortar Contact}

Mortar methods impose contact constraints in a weak, integral sense using
interpolation spaces for the contact pressure or multipliers. Instead of matching
the gap independently at isolated nodes, a weighted residual condition is used:

\begin{equation}
    \int_{\Gamma_c} N_{\lambda i} g_n \, d\Gamma = 0
\end{equation}

for the active contact region, where $ N_{\lambda i} $ are multiplier or mortar
shape functions.

Mortar methods are well suited to nonmatching meshes and can provide improved
conservation and reduced master-slave bias. Their formulation, integration and
active-set treatment are more involved than classical node-to-surface contact.

\subsection{One-Pass and Two-Pass Contact}

A one-pass method projects one selected side onto the other. It is inherently
asymmetric because exchanging master and slave changes the discrete equations.

A two-pass or symmetric method performs the constraint construction in both
directions. This may reduce bias, but simply duplicating constraints can also
overconstrain the interface, especially for matching meshes or nearly
incompressible response. A properly designed symmetric or mortar formulation is
not equivalent to blindly applying two independent node-to-surface pairs.

The words \textit{symmetric} and \textit{asymmetric} should therefore be read as
properties of the discrete contact algorithm, not of the physical friction law.

\section{Linearization and Nonlinear Solution}

The finite element equilibrium equations including contact may be written

\begin{equation}\label{contact-global-equilibrium}
    \m{r}(\m{u}) =
    \m{f}_{int}(\m{u})
    + \m{f}_c(\m{u})
    - \m{f}_{ext} = \m{0}.
\end{equation}

A Newton iteration solves

\begin{equation}\label{contact-newton-system}
    \m{K}_T^i \Delta \m{u}^i = -\m{r}^i,
    \qquad
    \m{u}^{i+1}=\m{u}^i+\Delta \m{u}^i,
\end{equation}

where

\begin{equation}
    \m{K}_T =
    \frac{\partial \m{f}_{int}}{\partial \m{u}}
    + \frac{\partial \m{f}_c}{\partial \m{u}}.
\end{equation}

The contact tangent contains more than a normal spring stiffness. Depending on
the formulation it includes derivatives of the gap, surface normal, projection
coordinates, pressure law and friction law. In finite sliding, changes of the
contact geometry produce additional geometric terms.

A typical contact iteration consists conceptually of:

\begin{enumerate}
    \item predict or update the current geometry,
    \item search for candidate surface pairs,
    \item compute projections, gaps and relative tangential increments,
    \item classify each contact point as open, closed-stick or closed-slip,
    \item assemble contact residuals and tangent contributions,
    \item solve the Newton correction,
    \item repeat until equilibrium and contact conditions are satisfied.
\end{enumerate}

The actual order may differ, and some solvers update status only at selected
stages to improve robustness.

\subsection{Contact Chattering}

A contact point may alternate between open and closed states in successive
iterations. Similarly, a friction point may alternate between stick and slip.
This behaviour is called \textit{chattering}. It can occur near the transition
where both branches are nearly admissible.

Chattering may be reduced by smaller load increments, line search, status
tolerances, augmented Lagrangian updates or mild regularization. Excessive
stabilization, however, can alter the physical force-displacement response.

\subsection{Initial Gap and Initial Penetration}

Before the first equilibrium iteration, the discretized surfaces may have an
initial gap or may overlap because of geometry tolerances, mesh approximation,
shell offsets or intentional interference.

An initial penetration may be treated in several conceptually different ways:

\begin{enumerate}
    \item preserve it and generate an initial contact force,
    \item shift the contact reference so that the initial configuration is
        considered just touching,
    \item move nodes or surfaces to remove the overlap,
    \item ramp the interference over several increments.
\end{enumerate}

These choices do not describe the same model. Removing an initial penetration
changes the reference geometry, while preserving it may introduce prestress.
The adopted treatment must therefore be understood rather than accepted as a
purely numerical setting.

\subsection{Search Distance and Candidate Detection}

To avoid testing every surface face against every other face, contact algorithms
use spatial search structures and a finite candidate distance. A contact pair
outside the search region is ignored during that search pass.

A search distance that is too small can miss rapidly closing or initially
separated surfaces. An excessively large distance increases search cost and can
create unintended candidate pairs. Search tolerance affects detection; it does
not by itself define the physical normal stiffness.

\section{Numerical Behaviour and Assessment}

Contact results should not be assessed only by whether the nonlinear solver
reports convergence. A converged solution may still contain excessive penalty
penetration, unintended contact pairs or mesh-dependent pressure peaks.

\subsection{Quantities to Check}

Useful checks include:

\begin{enumerate}
    \item maximum and representative normal penetration,
    \item normal contact pressure and integrated contact force,
    \item open, closed, sticking and sliding status distributions,
    \item frictional traction relative to the bound $ \mu p_n $,
    \item equilibrium residual and reaction-force balance,
    \item sensitivity to contact mesh refinement,
    \item sensitivity to penalty stiffness and load-step size.
\end{enumerate}

A penetration should be judged relative to a meaningful geometrical scale, such
as local element thickness, surface curvature radius or allowed interface
deflection. Reporting only its absolute value can be misleading when model units
or component sizes differ.

\subsection{Pressure Singularities}

High contact pressure near a sharp edge, corner or the end of a contact region
may be a genuine mathematical singularity. Mesh refinement can then increase
the peak pressure rather than converge it to a finite value. Integrated force,
contact width or pressure away from the singular point may be more meaningful
quantities.

A smooth pressure contour is not proof of correctness. Nodal averaging and
postprocessing can hide local oscillations or violations of the contact
conditions.

\subsection{Mesh Compatibility and Bias}

Nonmatching contact meshes are permissible, but the relative element sizes and
orders influence pressure transfer. A very coarse surface contacting a very fine
surface may produce oscillatory tractions or artificial local compliance.

Changing master and slave sides is a useful diagnostic for a biased one-pass
formulation. Large changes in global response indicate that the contact
discretization, not only the physical model, controls the answer.

\subsection{Contact in Dynamic Problems}

In dynamics, closing contact can introduce an impact. A purely displacement-based
penalty law stores elastic contact energy

\begin{equation}
    \Pi_c = \frac{1}{2}\epsilon_n \langle -g_n \rangle^2
\end{equation}

and may generate high-frequency oscillations. Contact damping or an impact law
may be added, but such additions dissipate energy and must be distinguished from
material damping.

The stable time step of an explicit method can be strongly reduced by a large
penalty stiffness because the contact stiffness increases the highest system
frequency. Thus the penalty value affects not only penetration but also
computational cost.

For implicit dynamics, abrupt status changes and frictional slip can make the
Newton iterations difficult. Time-step reduction often helps because it limits
the amount of contact-status change within one increment.

\section{Relation to Multifreedom Constraints}

The connection to Chapter 3 can now be stated precisely.

A bonded interface is essentially a set of displacement equality constraints:

\begin{equation}
    \m{A}\m{u}=\m{0}.
\end{equation}

A no-separation normal interface is a bilateral equality constraint on the
normal gap. A sticking tangential state is an equality constraint on incremental
tangential slip.

Ordinary normal contact is different because it combines an inequality with a
reaction-sign condition and complementarity:

\begin{equation}
    g_n \geq 0, \qquad p_n \geq 0, \qquad p_n g_n=0.
\end{equation}

Friction adds a second inequality and a second status decision:

\begin{equation}
    \norm{\m{t}_t} - \mu p_n \leq 0.
\end{equation}

Therefore contact may be regarded as a changing collection of MFC-like
constraints whose membership is unknown in advance. Penalty, multiplier and
augmented Lagrangian methods reappear, but they are embedded in geometry search,
active-set logic and constitutive stick-slip updates.

This viewpoint is useful because it separates three questions that are often
mixed together:

\begin{enumerate}
    \item \textit{What relative motion is physically allowed?}
        This determines frictionless, frictional, rough, no-separation or bonded
        behaviour.

    \item \textit{How is the constraint represented on the mesh?}
        This determines node-to-node, node-to-surface, surface-to-surface or
        mortar contact.

    \item \textit{How is the discrete constraint imposed?}
        This determines penalty, Lagrange multiplier, augmented Lagrangian or
        direct elimination treatment.
\end{enumerate}

A contact setting is properly understood only when all three questions can be
answered.

\section{Summary}

Contact is a unilateral interaction. The normal gap may remain positive with
zero pressure, or the gap may close and develop a compressive pressure. These
alternatives are expressed by

\begin{equation}
    g_n \geq 0, \qquad p_n \geq 0, \qquad p_n g_n=0.
\end{equation}

The contact type describes admissible relative motion and transmitted traction:

\begin{enumerate}
    \item frictionless contact allows separation and sliding and transmits only
        compression,
    \item frictional contact adds a pressure-dependent shear limit,
    \item rough contact prevents sliding while closed but still allows separation,
    \item no-separation contact prevents opening but allows tangential sliding,
    \item bonded contact prevents both opening and sliding and can transmit
        tension, compression and shear.
\end{enumerate}

Penalty, Lagrange multiplier and augmented Lagrangian methods describe how the
normal and tangential constraints are enforced numerically. Node-to-node,
node-to-surface, surface-to-surface and mortar methods describe how the
continuous interface is discretized.

No single contact definition or algorithm is uniformly best. The appropriate
choice depends on the expected motion, required force transfer, mesh relation,
constraint accuracy and nonlinear robustness. Contact settings should therefore
be selected from their mechanical meaning first and tuned numerically second.
